% !TEX program = xelatex
% Заготовка отзыва руководителя ВКР по форме инструкции ВКР 2026:
% 2б_Отзыв_ВКР_Руковд(УрФУ)+Соруковд(УрФУ).doc
% Сборка: ./compile.sh report/report.tex
\documentclass[a4paper,12pt]{extarticle}

% extarticle с опцией 12pt берёт классический \@xiipt (12 pt / 14,5 pt),
% а бланк набран шагом 13,8 bp (одинарный интервал Word): 13,8 bp = 13,852 pt.
\makeatletter
\renewcommand\normalsize{\@setfontsize\normalsize{12}{13.852}}
\makeatother

\usepackage{fontspec}
\setmainfont{Times New Roman}
\usepackage{polyglossia}
\setdefaultlanguage{russian}
\newfontfamily\cyrillicfont{Times New Roman}[Script=Cyrillic]
% Поля бланка отзыва: 25,4 / 11,05 / 19 / 9,5 мм (sectPr эталона).
\usepackage[left=25.4mm,right=11.05mm,top=19mm,bottom=9.5mm]{geometry}
\usepackage{array}

\pagestyle{empty}
\setlength{\parindent}{0pt}
\setlength{\parskip}{0pt}
% Первая базовая линия — на 11,34 bp ниже верхнего поля, как в эталоне.
\setlength{\topskip}{11.382pt}
\setlength{\arrayrulewidth}{0.5pt}
\tolerance=2000
\emergencystretch=3em

% Размеры бланка сняты с PDF-эталона в типографских пунктах (1 bp = 1/72",
% так считает Word) и переводятся в TeX pt множителем 803/800.
\newcommand{\vkrbp}[1]{\dimexpr#1pt*803/800\relax}
% Линия заполняемого поля: 0,5 pt, на 1,5 pt ниже базовой линии.
\newcommand{\vkrblank}[1]{\rule[-1.5pt]{\vkrbp{#1}}{0.5pt}}
\newcommand{\vkrline}[1]{\makebox[\linewidth][l]{#1}}
\newcommand{\vkrcenter}[1]{\makebox[\linewidth][c]{#1}}
% Абзац с распоркой: высота строки фиксирована, поэтому расстояние между
% базовыми линиями задаётся \vkrgap точно.
\newcommand{\vkrp}[1]{\strut #1\par}
% Расстояние между базовыми линиями соседних строк (в bp).
\newdimen\vkrskip
\newcommand{\vkrgap}[1]{%
  \par
  \ifdim\prevdepth>-100pt \vkrskip=\prevdepth \else \vkrskip=0pt \fi
  \prevdepth=-1000pt
  \vspace*{\dimexpr\vkrbp{#1}-9.6964pt-\vkrskip\relax}}
% Расстояние от низа предыдущего бокса (таблицы) до следующей базовой линии.
\newcommand{\vkrgapabs}[1]{%
  \par\prevdepth=-1000pt
  \vspace*{\dimexpr\vkrbp{#1}-9.6964pt\relax}}
% Распорка ячейки таблицы: высота и глубина строки как в бланке Word.
\newcommand{\vkrstrut}{\vrule height 11.2pt depth 2.85pt width 0pt}
% Строка таблицы в бланке — ровно 28,1 bp: подрезаем пустое место под текстом.
\newcommand{\vkrrowfix}{-1.505pt}

% Пустой макрос оставляет ветку бланка. Сравнение — по уже раскрытому
% значению: \newcommand делает макрос \long, поэтому \ifx с \@empty не годится.
\makeatletter
\newcommand{\vkrifempty}[1]{%
  \expandafter\vkr@ifempty\expandafter{#1}}
\newcommand{\vkr@ifempty}[1]{%
  \if\relax\detokenize{#1}\relax
    \expandafter\@firstoftwo
  \else
    \expandafter\@secondoftwo
  \fi}
\newbox\vkrbox
\newbox\vkrlb
\newbox\vkrlines
\newdimen\vkrprefixwd
\newcount\vkrleft
\newcount\vkrn
\makeatother
% Ширина полной линии бланка. Поле на ней длиннее текста: линия идёт
% до этого края, а не обрывается на последней букве.
\newcommand{\vkrfieldwd}{\vkrbp{486}}
% Поле фиксированной ширины (в bp). Пустой текст — линия бланка. Текст
% короче поля лежит на линии заданной ширины. Более длинный остаётся
% на одной строке, линия на 18 bp длиннее текста.
\newcommand{\vkrfield}[2]{%
  \expandafter\vkrfieldbody\expandafter{#2}{#1}}
\newcommand{\vkrfieldbody}[2]{%
  \if\relax\detokenize{#1}\relax
    \vkrblank{#2}%
  \else
    \begingroup
    \setbox0=\hbox{#1}%
    \ifdim\wd0>\vkrbp{#2}\relax
      \makebox[\dimexpr\wd0+\vkrbp{18}\relax][l]{%
        \rlap{\rule[-1.5pt]{\dimexpr\wd0+\vkrbp{18}\relax}{0.5pt}}#1}%
    \else
      \makebox[\vkrbp{#2}][l]{\rlap{\vkrblank{#2}}#1}%
    \fi
    \endgroup
  \fi}
% Строка «метка + значение». Линия начинается у значения и доходит до
% правого края поля, поэтому она длиннее текста.
\newcommand{\vkrvalueline}[2]{%
  \begingroup
  \setbox0=\hbox{#1}%
  \setbox1=\hbox{#2}%
  \vkrprefixwd=\dimexpr\vkrfieldwd-\wd0\relax
  \ifdim\wd1>\vkrprefixwd
    \vkrprefixwd=\dimexpr\wd1+\vkrbp{18}\relax
  \fi
  \vkrp{\hbox to \linewidth{\box0\rlap{\rule[-1.5pt]{\vkrprefixwd}{0.5pt}}\box1\hfil}}%
  \endgroup}
% Снять клей и штраф над последним боксом: между строками абзаца сначала
% идёт клей, и один \unpenalty его не достаёт.
\newcommand{\vkrstrip}{%
  \ifnum\lastnodetype=11
    \unskip\vkrstrip
  \else\ifnum\lastnodetype=12
    \unkern\vkrstrip
  \else\ifnum\lastnodetype=13
    \unpenalty\vkrstrip
  \fi\fi\fi}
% Собрать строки абзаца снизу вверх и провести поле на каждой.
\newcommand{\vkrrepack}{%
  \vkrstrip
  \setbox0=\lastbox
  \ifvoid0
  \else
    \begingroup
    \setbox\vkrlb=\box0
    \vkrrepack
    \nointerlineskip
    \hbox to \vkrfieldwd{%
      \ifnum\vkrn=0
        \rlap{\hspace*{\vkrprefixwd}\rule[-1.5pt]{\dimexpr\vkrfieldwd-\vkrprefixwd\relax}{0.5pt}}%
      \else
        \rlap{\rule[-1.5pt]{\vkrfieldwd}{0.5pt}}%
      \fi
      \box\vkrlb\hfil}%
    \global\advance\vkrn by 1
    \endgroup
  \fi}
% Многострочное поле бланка. #1 — шаг базовых линий (bp), #2 — число строк
% вместе со строкой метки, #3 — метка, #4 — текст. Текст кончается на 18 bp
% раньше линии, поэтому линия длиннее текста. Строки выводятся по одной
% через \vkrgap: пока текст умещается в #2 строк, шаг совпадает с пустым блоком.
\newcommand{\vkrtextbox}[4]{%
  \begingroup
  \settowidth{\vkrprefixwd}{#3}%
  \setbox\vkrbox=\vbox{%
    \hsize=\vkrfieldwd
    \parindent=0pt
    \parskip=0pt
    \baselineskip=\vkrbp{#1}\relax
    \lineskiplimit=-100pt
    \raggedright
    \setlength{\rightskip}{\vkrbp{18}}%
    \addtolength{\rightskip}{0pt plus 1fil}%
    \noindent\strut #3#4\par
    \global\vkrleft=\prevgraf
  }%
  \ifnum\vkrleft>#2\relax
    \PackageWarning{vkrreport}{Текст поля занял \the\vkrleft\ строк вместо #2}%
  \fi
  \vkrn=0
  \global\setbox\vkrbox=\vbox{%
    \unvbox\vkrbox
    \vkrrepack
    \loop
    \ifnum\vkrn<#2\relax
      \nointerlineskip
      \hbox to \vkrfieldwd{\strut\rlap{\rule[-1.5pt]{\vkrfieldwd}{0.5pt}}\hfil}%
      \global\advance\vkrn by 1
    \repeat
  }%
  \gdef\vkrstep{#1}%
  \endgroup
  \global\setbox\vkrlines=\vbox{}%
  \loop
    \setbox0=\vbox{\unvbox\vkrbox\vkrstrip\global\setbox\vkrlb=\lastbox}%
    \global\setbox\vkrbox=\box0
  \ifvoid\vkrlb
  \else
    \global\setbox\vkrlines=\vbox{\unvbox\vkrlines\nointerlineskip\box\vkrlb}%
  \repeat
  \vkrn=0
  \loop
    \setbox0=\vbox{\unvbox\vkrlines\vkrstrip\global\setbox\vkrlb=\lastbox}%
    \global\setbox\vkrlines=\box0
  \ifvoid\vkrlb
  \else
    \ifnum\vkrn>0
      \vkrgap{\vkrstep}%
    \fi
    \vkrp{\vkrline{\box\vkrlb}}%
    \advance\vkrn by 1
  \repeat}
% Направление на строке бланка. Пустые код и название ничего не добавляют.
\newcommand{\vkrdirectionvalue}{%
  \vkrifempty{\VKRDirectionCode}{%
    «\VKRDirection»%
  }{\VKRDirectionCode\vkrifempty{\VKRDirection}{}{ «\VKRDirection»}}%
}
% В заключении пустая пара код+название оставляет подсказку бланка.
\newcommand{\vkrdirectionconclusion}{%
  \vkrifempty{\VKRDirectionCode}{%
    \vkrifempty{\VKRDirection}{%
      \textbf{код \mbox{направления «название направления»}}%
    }{%
      \textbf{\mbox{«\VKRDirection»}}%
    }%
  }{%
    \textbf{\mbox{\VKRDirectionCode\vkrifempty{\VKRDirection}{}{ «\VKRDirection»}}}%
  }%
}

% —————————————————————————
% Пример заполнения. Студент, код и название направления, программа
% и руководитель согласованы с thesis/thesis.tex. Тема та же, обычным
% написанием: в ВКР она хранится прописными для титула. Пустое значение {}
% возвращает линию или подсказку бланка.
% \VKRDirection — наименование без кавычек: бланк сам ставит «».
% \VKRAdvisor и \VKRSorukovoditel — «Ф. И. О., учёная степень, учёное звание
% (должность)». Заполненные направление, программа, руководитель и
% соруководитель подчёркиваются. Тема, достоинства, недостатки и
% информация для ГЭК переносятся по строкам бланка (2, 4, 3 и 3 строки).
% —————————————————————————
\newcommand{\VKRTitle}{Хроматографическое исследование системы полиэтиленоксид–хитозан}
\newcommand{\VKRStudent}{Смирнов Юрий Викторович}
\newcommand{\VKRGruppa}{МЕН-440301}
\newcommand{\VKRDirectionCode}{04.03.01}
\newcommand{\VKRDirection}{Химия}
\newcommand{\VKRProgram}{«Химия»}
\newcommand{\VKRAdvisor}{Иванов Иван Сергеевич, к. х. н., доц.}
\newcommand{\VKRSorukovoditel}{Петрова Анна Сергеевна, к. х. н., доц.}
\newcommand{\VKRMerits}{Студент самостоятельно освоил хроматографический анализ системы полиэтиленоксид–хитозан, аккуратно обработал результаты и ясно их изложил. Выводы согласуются с полученными данными.}
\newcommand{\VKRShortcomings}{Недостаточно подробно оценена погрешность измерений. Сопоставление с литературными данными можно расширить.}
\newcommand{\VKRGEKInfo}{Работу рекомендуется допустить к защите. Результаты пригодны для учебного практикума по хроматографии полимеров.}
\newcommand{\VKRQualification}{бакалавр}

\begin{document}

\vkrp{\vkrcenter{Министерство науки и высшего образования Российской Федерации}}
\vkrgap{13.8}
\vkrp{\vkrcenter{Федеральное государственное автономное образовательное учреждение высшего образования}}
\vkrgap{13.8}
\vkrp{\vkrcenter{«Уральский федеральный университет имени первого Президента России Б.Н. Ельцина»}}

\vkrgap{34.5}
\vkrp{\vkrcenter{\textbf{ОТЗЫВ}}}
\vkrgap{20.7}
\vkrp{\vkrcenter{руководителя выпускной квалификационной работы}}
\vkrgap{20.7}
\vkrifempty{\VKRTitle}{%
\vkrp{\vkrline{Выпускная квалификационная работа выполнена на тему\hspace{3.1pt}\vkrblank{186}}}
\vkrgap{20.7}
\vkrp{\vkrline{\vkrblank{486}}}
}{%
\vkrtextbox{20.7}{2}{Выпускная квалификационная работа выполнена на тему\hspace{3.1pt}}{\VKRTitle}%
}
\vkrgap{20.7}
\vkrp{\vkrline{Студентом\hspace{3.1pt}\vkrfield{264}{\VKRStudent}\hspace{3pt}группы\hspace{3.1pt}\vkrfield{120}{\VKRGruppa}}}
\vkrgap{20.7}
\vkrifempty{\VKRDirectionCode}{%
  \vkrifempty{\VKRDirection}{%
    \vkrp{Направление (специальность):}%
  }{%
    \vkrvalueline{Направление (специальность): }{\vkrdirectionvalue}%
  }%
}{%
  \vkrvalueline{Направление (специальность): }{\vkrdirectionvalue}%
}
\vkrgap{13.8}
\vkrifempty{\VKRProgram}{%
  \vkrp{Образовательная программа:}%
}{%
  \vkrvalueline{Образовательная программа: }{\VKRProgram}%
}
\vkrgap{13.8}
\vkrifempty{\VKRAdvisor}{%
  \vkrp{Руководитель: Ф.И.О., ученая степень, ученое звание (должность)}%
}{%
  \vkrvalueline{Руководитель: }{\VKRAdvisor}%
}
\vkrgap{13.8}
\vkrifempty{\VKRSorukovoditel}{%
  \vkrp{Соруководитель: Ф.И.О., ученая степень, ученое звание (должность)}%
}{%
  \vkrvalueline{Соруководитель: }{\VKRSorukovoditel}%
}

\vkrgap{19.5}
% В бланке у этого абзаца выравнивание по левому краю (jc не задан).
\begingroup\raggedright
\vkrp{Оценка соответствия требованиям ФГОС ВО подготовленной выпускной квалификационной работы}
\endgroup

\vkrgap{12.26}
\newsavebox{\vkrtable}
\sbox{\vkrtable}{%
\begingroup
\setlength{\tabcolsep}{5.4pt}%
\setlength{\extrarowheight}{0.25pt}%
\begin{tabular}{|p{159.35pt}|>{\centering\arraybackslash}p{74pt}|>{\centering\arraybackslash}p{74.81pt}|>{\centering\arraybackslash}p{72.17pt}|>{\centering\arraybackslash}p{56.46pt}|}
\hline
\vkrstrut Результаты освоения основной профессиональной образовательной программы, представленные в ВКР
  & \vkrstrut Полностью соответствует
  & \vkrstrut Соответствует
  & \vkrstrut В основном соответствует
  & \vkrstrut Не соответствует \\[\vkrrowfix]
\hline
\vkrstrut В части общекультурных компетенций & \vkrstrut & \vkrstrut & \vkrstrut & \vkrstrut \\[\vkrrowfix]
\hline
\vkrstrut В части общепрофессиональных компетенций & \vkrstrut & \vkrstrut & \vkrstrut & \vkrstrut \\[\vkrrowfix]
\hline
\vkrstrut В части профессиональных компетенций & \vkrstrut & \vkrstrut & \vkrstrut & \vkrstrut \\[\vkrrowfix]
\hline
\vkrstrut В части дополнительных компетенций (при наличии) & \vkrstrut & \vkrstrut & \vkrstrut & \vkrstrut \\[\vkrrowfix]
\hline
\end{tabular}%
\endgroup}
\noindent\usebox{\vkrtable}\par

\vkrgapabs{20.42}
\vkrifempty{\VKRMerits}{%
\vkrp{\vkrline{\textbf{Отмеченные достоинства:}\hspace{3pt}\vkrblank{336}}}
\vkrgap{13.8}
\vkrp{\vkrline{\vkrblank{486}}}
\vkrgap{13.8}
\vkrp{\vkrline{\vkrblank{486}}}
\vkrgap{13.8}
\vkrp{\vkrline{\vkrblank{486}}}
}{%
\vkrtextbox{13.8}{4}{\textbf{Отмеченные достоинства:}\hspace{3pt}}{\VKRMerits}%
}
\vkrgap{13.8}
\vkrifempty{\VKRShortcomings}{%
\vkrp{\vkrline{\textbf{Отмеченные недостатки:}\hspace{3pt}\vkrblank{342}}}
\vkrgap{13.8}
\vkrp{\vkrline{\vkrblank{486}}}
\vkrgap{13.8}
\vkrp{\vkrline{\vkrblank{486}}}
}{%
\vkrtextbox{13.8}{3}{\textbf{Отмеченные недостатки:}\hspace{3pt}}{\VKRShortcomings}%
}

\vkrgap{27.6}
\vkrifempty{\VKRGEKInfo}{%
\vkrp{\vkrline{Дополнительная информация для ГЭК:\hspace{3pt}\vkrblank{282}}}
\vkrgap{13.8}
\vkrp{\vkrline{\vkrblank{486}}}
\vkrgap{13.8}
\vkrp{\vkrline{\vkrblank{486}}}
}{%
\vkrtextbox{13.8}{3}{Дополнительная информация для ГЭК:\hspace{3pt}}{\VKRGEKInfo}%
}

\vkrgap{27.6}
\vkrp{\textbf{Заключение.} Выполненная работа является самостоятельным, законченным
исследованием,
\mbox{достойным/недостойным} представления перед ГЭК. Студент (Ф.И.О.)
\vkrfield{108}{\VKRStudent} заслуживает/не заслуживает присуждения квалификации
\vkrifempty{\VKRQualification}{\textbf{бакалавр}}{\textbf{\mbox{\VKRQualification}}} по направлению \vkrdirectionconclusion.}

\vkrgap{41.4}
\vkrp{\vkrline{Руководитель:\hspace{3pt}\vkrblank{165}\hspace{27.1pt}«\hspace{3pt}\vkrblank{24}\hspace{3pt}»\hspace{6pt}\vkrblank{78}\hspace{6pt}20\vkrblank{24}\hspace{3pt}г.}}
\vkrgap{13.8}
\vkrp{\vkrline{\hspace*{135pt}(подпись)}}

\vkrgap{27.6}
\vkrp{\vkrline{Соруководитель:\hspace{3pt}\vkrblank{165}\hspace{27.1pt}«\hspace{3pt}\vkrblank{24}\hspace{3pt}»\hspace{6pt}\vkrblank{78}\hspace{6pt}20\vkrblank{24}\hspace{3pt}г.}}
\vkrgap{13.8}
\vkrp{\vkrline{\hspace*{135pt}(подпись)}}

\end{document}
